\begin{tikzpicture}[paperfig]
\node[anchor=east] at(-.2,1.45){$K=4,\ j=1$};
\draw[line width=1pt] (0,1.45)--(9,1.45);
\foreach \x/\l in {0/{1/4},9/{1/2}} {\node[above=3pt] at(\x,1.45){$\l$};}
\fill (0,1.45) circle(2pt);\draw[fill=white] (9,1.45) circle(2pt);
\node[anchor=east] at(-.2,0){$fK=12$};
\foreach \x/\j in {0/3,3/4,6/5}{
 \draw[line width=1pt] (\x,0)--++(3,0);
 \fill (\x,0) circle(2pt);\draw[fill=white] (\x+3,0) circle(2pt);
 \node[below=3pt] at(\x+1.5,0){$\j$};
}
\node[above=3pt] at(3,0){$1/3$};\node[above=3pt] at(6,0){$5/12$};
\draw[fedge,dashed] (4.5,1.2)--(4.5,.3) node[midway,left] {$f=3$};
\draw[femph,dotted] (5.4,1.45)--(5.4,0);
\fill[figblue] (5.4,1.45) circle(2pt);\fill[figblue] (5.4,0) circle(2pt);
\node[above=3pt,align=center] at(5.4,1.5){$\theta=2/5$};
\node[fbox,text width=11.5cm] at (3.8,-1.25)
 {$\lfloor4\theta\rfloor=1,\qquad\lfloor12\theta\rfloor=4=3\cdot1+1$\\[5pt]
 $\sigma_K(n)\in U\quad\Longrightarrow\quad
 \sigma_{fK}(n)\in\operatorname{Refine}_f(U)$};
\node[align=center,text width=12cm] at(3.8,-2.5)
 {\figtext{The same true orbit point lies in exactly one finer cell.}{同じ真の軌道点は、細分後のちょうど一つのセルに入る。}\\
 \figtext{Edges are rebuilt; an arbitrary coarse walk need not lift.}{辺は再構成する。任意の粗いウォークが持ち上がるとは限らない。}};
\end{tikzpicture}
